\section{Overview of Risk Management Tools}
Risk management tools are used to support and structure the identification, analysis, evaluation, and treatment of information security risks. In practice, such tools vary widely in scope, complexity, and intended use. Some focus primarily on documenting risks and controls, while others provide broader support for governance, compliance, and integration with existing security and IT management processes. The level of automation and methodological support also differs significantly between tools.

\subsection{Tool categories}
In order to cover a broad and representative selection of risk management tools, this thesis includes both open-source and commercial solutions commonly used within information security risk management. The distinction between these categories was chosen to reflect differences in licensing models, availability, and typical adoption contexts, which are relevant factors when assessing practical applicability.

\subsubsection{Open-source tools}
Open-source risk management tools are typically developed by communities, research institutions, or non-commercial organizations. They often allow organizations to inspect, modify, and adapt the tool to their specific needs. Such tools may be particularly attractive to organizations with limited budgets or strong technical expertise, but they can also require more effort in terms of setup, maintenance, and customization. \\

We have selected the following open-source tools for our research:

\begin{itemize}
    \item OpenRMF OSS
    \item Risk BowTie
    \item SimpleRisk Core
    \item tool 4    % TODO
    \item tool 5    % TODO
\end{itemize}

\subsubsection{Commercial tools}
Commercial risk management tools are generally offered as licensed products or subscription-based services. These tools often provide more extensive feature sets, structured workflows, and professional support. They are commonly positioned as part of broader governance, risk, and compliance (GRC) platforms and may offer tighter integration with enterprise systems and reporting requirements. \\

We have selected the following commercial tools for our research:

\begin{itemize}
    \item tool 1    % TODO
    \item tool 2    % TODO
    \item tool 3    % TODO
    \item tool 4    % TODO
\end{itemize}
